\section{El perfil canónico discreto}\label{sec:perfil}

El paso del tablero al cuadrado unidad exige comparar magnitudes de la misma escala.
Como las $n$ casillas de cada fila suman uno, su peso medio es $1/n$. Por ello comparamos
$n\pi_{ij}$, y no el peso sin reescalar, con el valor de $\nu_1$ en el centro de la casilla.

Primero hay que asegurar que el maximizador se mantiene alejado de los extremos en el sentido
preciso de la Sección~\ref{sec:P02}: sus pesos están entre $c/n$ y $M/n$ y las vacancias
diagonales tienen una cota inferior positiva. La Sección~\ref{sec:P01} prepara esas cotas
mediante un punto de Slater, es decir, un punto factible que cumple con desigualdad estricta las
restricciones diagonales y de casilla.
Con ellas se controla el espectro de $S_n$ en la Sección~\ref{sec:P03} y se desarrolla el
término local en la Sección~\ref{sec:P04}. Finalmente, las Secciones~\ref{sec:P05}
y~\ref{sec:P06} establecen la comparación anunciada con error uniforme $O(1/n)$.
